\naslov{Osnove Newtonove mehanike}

\begin{definicija}
  \pojem{Afin prostor} $\affin$ nad vektorskim prostorom $V$ je množica z
  binarno operacijo $+: \affin \times V \to \affin$, za katero velja:
  \begin{itemize}
  \item Za poljuben $A \in \affin$ ter $\mb{a},\mb{b} \in V$ velja $(A +
    \mb{a}) + \mb{b} = A + (\mb{a} + \mb{b})$
  \item Za poljubna $A, B \in \affin$ obstaja natanko določen $\mb{a} \in
    V$, da je $B = A + \mb{a}$.
  \end{itemize}
  \pojem{Dimenzija} afinega prostora je enaka dimenziji vektorskega prostora
  $V$.
\end{definicija}

\begin{definicija}
  Naj bo $\affin$ afin prostor nad vektorskih prostorom $V$.
  Definiramo operacijo odštevanja $\affin \times \affin \to V$ s
  predpisom
  \[
    B - A = \mb{a} \nt B = A + \mb{a}.
  \]
\end{definicija}

\begin{trditev}
  V afinem prostoru veljajo naslednje zveze:
  \begin{itemize}
  \item $A - A = \mb{0}$.
  \item $(A - B) + (B - A) = \mb{0}$.
  \item $(A - B) + (B - C) + (C - A) = \mb{0}$.
  \item $(A - B) + \mb{a} = (A + \mb{a}) - B$.
  \item $(A - B) + C = (C - B) + A$.
  \end{itemize}
\end{trditev}

\begin{definicija}
  Preslikava $g : \affin \to \affin'$ med afinima prostoroma je
  \pojem{afina}, če obstaja $dg \in L(V, V')$, da za vsaka $A, B \in
  \affin$ velja $g(A) - g(B) = dg(A - B)$.
\end{definicija}

Za afino preslikavo $g$ si lahko izberemo \pojem{pol} $O$, ter izpeljemo
\[
  g(A) = g(O) + dg(A - O).
\]
Vrednosti funkcije seveda niso odvisne od izbire pola.

\vprasanje{Definiraj afin prostor in afino preslikavo.}

\newcommand{\cas}{\mathfrak{t}}
\begin{definicija}
  \pojem{Galilejeva struktura} je trojica $\mathcal{G} = (\affin,
  \cas, \rho)$, kjer je $\affin$ štirirazsežni afin prostor nad $V$, $\cas \in
  L(V, \R)$ in $\rho$ ekvlidska metrika na $\jedro \cas$, porojena z normo
  $\norm{\cdot}$.
  Funkciji $\cas$ pravimo \pojem{časovnost}, elementom $\affin$ pa pravimo
  \pojem{dogodki}.
  Pretečeni čas med dogodkoma $A$ in $B$ označimo s $\cas(A, B)$.
  Dogodka sta \pojem{istočasna}, če je $\cas(A, B) = 0$.
  Za istočasne dogodke lahko definiramo razdaljo $\rho(A, B) = \norm{B - A}$
  (uporabimo isto oznako kot za metriko v $\jedro\cas$).
\end{definicija}

\begin{definicija}
  Galilejevi strukturi $\mathcal{G} = (\affin, \cas, \rho)$ in $\mathcal{G}' =
  (\affin', \cas', \rho')$ sta \pojem{ekvivalentni}, če obstaja afina bijekcija
  $g: \affin \to \affin'$, ki ohranja časovnost in razdaljo med istočasnimi
  dogodki;
  \begin{align*}
    \cas'(g(A) - g(B)) = \cas(A, B),
    &&
       \rho'(g(A), g(B)) = \rho(A, B).
  \end{align*}
  Taki transformaciji pravimo \pojem{Galilejeva transformacija}.
\end{definicija}

\vprasanje{Definiraj Galilejevo strukturo in Galilejeve transformacije.}

\newcommand{\galileo}{\R\times\mathbb{E}}
Modelni primer je naravna Galilejeva struktura na $\affin = \galileo$, kjer je
$\mathbb{E}$ trirazsežni Evklidski prostor.
Za elemente $A_i = (t_i, \mb{P}_i) \in \affin$ naravne strukture velja
\begin{itemize}
\item $\cas(A_1 - A_2) = t_1 - t_2$,
\item $\rho(A_1, A_2) = \norm{\mb{P}_1 - \mb{P}_2}$.
\end{itemize}

\begin{definicija}
  \pojem{Koordinatni sistem} na $\affin$ je bijekcija $\phi : \affin \to
  \galileo$ s komponentami $\phi(A) = (\tau \phi(A), \pi \phi(A))$, in pri
  kateri je $\tau \circ \phi$ linearna preslikava.
\end{definicija}

\begin{opomba}
  Če sta $\phi$ in $\phi'$ koordinatna sistema, je preslikava $\phi' \circ
  \phi^{-1}: \galileo \to \galileo$ bijekcija.
\end{opomba}

\vprasanje{Kaj je koordinantni sistem?}

\begin{izrek}
  Galilejeva transformacija $g: \galileo \to \galileo$ je oblike
  \[
    g(t, \mb{P}) = (t_0' + t, \mb{P}_0' + \vec{c} t + \mb{Q}(\mb{P} -
    \mb{P}_0)),
  \]
  kjer je $\mb{Q} \in O(3)$ ortogonalna transformacija.
\end{izrek}

\begin{proof}
  Ker je $g$ afina preslikava, jo lahko zapišemo kot
  \[
    g(t, \mb{P}) = g(t_0, \mb{P}_0) + dg(t - t_0, \mb{P} - \mb{P}_0),
  \]
  kjer je $dg \in L(\R^4, \R^4)$.
  Če označimo $g(t_0, \mb{P}_0) = (t_0', \mb{P}_0')$, in zapišemo $dg$ kot
  bločno matriko, dobimo
  \[
    g(t, \mb{P})
    =
    (t_0', \mb{P}_0')
    +
    \begin{bmatrix}
      \alpha & \vec{a}^T \\
      \vec{c} & \mb{Q}
    \end{bmatrix}
    \begin{bmatrix}
      t - t_0 \\ \mb{P} - \mb{P}_0
    \end{bmatrix}
    =
    (t_0', \mb{P}_0')
    +
    \begin{bmatrix}
      \alpha (t - t_0) + \vec{a} \cdot (\mb{P} - \mb{P}_0) \\
      (t - t_0).\vec{c} + \mb{Q} (\mb{P} - \mb{P}_0)
    \end{bmatrix}.
  \]
  Za dogodka $(t_1, \mb{P}_1)$ in $(t_2, \mb{P}_2)$ zahtevamo
  \[
    t_2 - t_1 = \tau(g(t_2, \mb{P}_2) - g(t_1, \mb{P}_1)).
  \]
  Če razvijemo desno stran zahteve po izpeljani formuli, dobimo pogoj
  \[
    t_2 - t_1 = \alpha (t_2 - t_1) + \vec{a} \cdot (\mb{P}_2 - \mb{P}_1).
  \]
  Iz tega sledi $\alpha = 1$ in $\vec{a} = \vec{0}$.
  Drug pogoj je, da se mora razdalja med istočasnimi dogodki ohranjati.
  Iz spodnjega dela bločne matrike dobimo pogoj
  \[
    \norm{\mb{P}_2 - \mb{P}_1} = \norm{ \mb{Q} (\mb{P}_2 - \mb{P}_1)},
  \]
  torej mora biti $\mb{Q}$ ortogonalna.
\end{proof}

\vprasanje{Kakšno obliko imajo Galilejeve transformacije $\galileo \to
  \galileo$? Dokaži.}